\section{What a quasipolynomial completion must prove}
\label{sec:global-accounting}

The preceding operations address different terms in a hierarchy's cost.
The genus budget limits the network joining retained pieces.  The native
moment compiler limits the description of a boundary response.  Unmarked
reconstruction limits the size of one new triangulation when the cutting
surface has controlled topological complexity.  None of these statements
bounds the number of hierarchy states visited, or the cost of discovering
the next useful surface.  This section makes that remaining distinction
quantitative.

\subsection{Four costs that must be bounded together}

For a state $H$, write $s(H)$ for its complete encoded size, including its
manifold presentation, boundary pattern, normal-coordinate charts,
compressed component multiplicities, and any proof data the algorithm
must retain.  This is deliberately more than the tetrahedron count.  A
prospective algorithm must bound the following quantities.

\begin{enumerate}
\item \emph{Search volume:} the number $J(n)$ of state visits or elementary
hierarchy iterations processed on an $n$-crossing input.  A repeated visit
is counted again.
\item \emph{State volume:} the maximum value $S(n)=\max_H s(H)$ during
that computation, including the largest binary integer lengths.
\item \emph{Local work:} the work needed to discover, perform and verify
the permitted transition from a state of size $s$.  A polynomial-time
checker for a supplied surface does not bound its discovery cost.
\item \emph{Completion:} a proof that every permitted branch is handled,
that the selected hierarchy operations remain sound, and that the final
test distinguishes all knots.  A resource cap returning
\code{INCONCLUSIVE} is an honest software contract, but not completeness.
\end{enumerate}

If local work is at most $S(n)^a$ for a fixed constant $a$, then the direct
accounting bound is $J(n)S(n)^a$, plus input construction.  This elementary
product is useful because it exposes exactly where a claimed fast local
primitive can and cannot change the overall exponent.

\begin{theorem}[Conditional hierarchy accounting]
\label{thm:conditional-accounting}
Suppose a complete and correct recognition procedure on diagrams with
$n\geq2$ crossings has parameters $L=L(n)\geq1$ and $G=G(n)\geq1$.
Assume its input construction takes polynomial time in $n$, and that all
subsequent work, including its final decision, is charged to the iterations
below.
Assume fixed constants $a,c,C>0$ and $A>1$ satisfy all of the following.
\begin{enumerate}[label=(\roman*)]
\item It performs at most $C(L+1)(G+1)^L$ elementary iterations, counting
revisits to a state again.
\item $G\leq n^c$, and every encoded state and transition output has
size at most $n^c A^{cL}(L+1)^c$.
\item Discovery, source validation, exact local computation, certificate
production and replay for one iteration together cost at most the $a$th
power of this size, times a fixed constant.
\end{enumerate}
Then its total running time is
\[
 \exp\!\bigl(O((L+1)\log n+L+\log(L+1))\bigr).
\]
In particular, $L=O((\log n)^2)$ gives
$2^{O((\log n)^3)}$, while $L=O(\log n)$ gives
$n^{O(\log n)}$.
\end{theorem}
\begin{proof}
Multiply the bounds in (i) and (iii), using (ii), and take logarithms.
Up to constant factors, the logarithm of the iteration cost is at most
\[
 \log(L+1)+L\log(G+1)
 +ac\log n+acL\log A+ac\log(L+1).
\]
Since $\log(G+1)=O(\log n)$ and $A,a,c$ are fixed, this is the stated
expression.  Adding the polynomial input-construction cost does not
change it, since that cost has logarithm $O(\log n)$.  Substituting either
bound on $L$ proves the last two claims.
\end{proof}

This is an implication, not a new hierarchy-existence theorem.  Its first
assumption matches the form of the iteration bound in Proposition~9.1 of
the 2026 preprint \cite{Lackenby2026Hierarchies}; a finite search tree of
depth $L$ and branching at most $G+1$ also has the stated coarse bound.
The geometric meanings of $L$ and $G$ must come from the actual procedure,
not from arbitrary tuning parameters.  The theorem explains why the two
2021 depth announcements lead to genuinely different exponents.

\subsection{What the unmarked reconstruction bound contributes}

\begin{proposition}[Iteration of an unmarked size bound]
\label{prop:reset-recurrence}
Suppose a sequence of unmarked reconstructions satisfies
\[
 t_{i+1}\leq A(t_i+k_i+1),\qquad A\geq1,
\]
where $k_i=b_0(S_i)+\chi_-(S_i)$ counts every cutting component.
Then
\[
 t_i\leq A^i\left(t_0+\sum_{j<i}(k_j+1)\right).
\]
If $t_0,k_j\leq n^c$ and $i\leq L$, then
\[
 \log(t_i+1)=O(L+\log(L+1)+\log n)
\]
for fixed $A,c$.
\end{proposition}
\begin{proof}
For $i=0$ the assertion is equality.  Given the bound for $i$, multiply
by $A$ and use $A\leq A^{i+1}$ to absorb the next additive term.  This
proves the first inequality by induction.  The sum is at most
$i(n^c+1)$, which gives the logarithmic estimate.
\end{proof}

Thus even a constant-factor expansion at each of
$O((\log n)^2)$ cuts can remain compatible with the later announced
quasipolynomial exponent.  A perfectly size-preserving reset is not
necessary for this accounting.  What is necessary is control of the
\emph{whole} state: the statement above does not bound the word describing
a carried boundary pattern, the normal coordinates of the next surface,
or the search required to find it.  If the remainder of each state has
\emph{total} encoded size polynomial in $t_i$ and $n$, the encoded-state
part of assumption (ii) in Theorem~\ref{thm:conditional-accounting} follows
after adjusting the fixed constants.  A polynomial bound for $G$ and the
discovery-time bound in (iii) remain separate obligations.  Small output
size alone supplies neither of them.

Nor can one discard multiplicities to make $k_i$ appear small.  The
nested-sphere example in Section~\ref{sec:unmarked-reconstruction} shows
that a single type record can represent exponentially many components.
A hierarchy may have a valid rule for purging those components; otherwise
their explicit reconstruction is a real output cost.  Any proposed reset
must state which of these two contracts it implements.

\subsection{Sharing already supplied transport along a forest}

The rooted forest of Section~\ref{sec:attachment-kernel} gives a precise
place to store transport information without expanding every source route.
There is an elementary algebraic principle behind this design.

\begin{proposition}[Forest transport circuit]
\label{prop:forest-transport}
Let $Q$ be a finite oriented graph with a rooted spanning forest $F$, and
suppose every oriented edge $e$ has a supplied label $a_e$ in a group,
with $a_{\bar e}=a_e^{-1}$.  Labels may instead be words in a formal group
alphabet, represented by shared expression nodes.  There is a circuit with
$O(|V(Q)|+|E(Q)|)$ additional product and inverse gates that labels every
nonforest edge of the contracted graph and preserves the products around
closed walks based at roots.  Only the original label descriptions are
needed as input leaves.

If group multiplication and inversion have unit cost, all contracted
labels can be evaluated using $O(|V(Q)|+|E(Q)|)$ group operations.  The
claim is about supplied algebraic transport; it does not construct a
geometric transport map on a boundary pattern.
\end{proposition}
\begin{proof}
Let $P_v$ be the product along the forest route from $v$ to its root,
with products written in path order.  For a root set $P_v=1$; otherwise
define a shared gate
\[
 P_v=a_{(v,\operatorname{parent}(v))}\,P_{\operatorname{parent}(v)}.
\]
For an edge $e:u\longrightarrow v$ not in the forest, assign
\[
                   a'_e=P_u^{-1}a_eP_v.
\]
Each vertex requires a bounded number of gates and each retained edge
requires a bounded number more.  A forest edge has transformed label one.
Along any walk, the consecutive $P_vP_v^{-1}$ factors cancel, leaving
$P_{\rm start}^{-1}$, the original path product, and $P_{\rm end}$.
For a closed walk based at a root both outside factors are one.  Processing
the forest from roots to leaves and then processing the nonforest edges
gives the operation count.  No path is expanded as a word.
\end{proof}

This supplies a useful interface for a future patterned reducer.  Local
fibre signs can be evaluated in $C_2$ immediately.  Longer curve or chart
instructions can remain as a circuit until a consumer requests them.
The missing step is geometric: prove that the proposed source labels
encode the required maps, provide compatible domains, and bound the cost
of the subsequent equality or isotopy problem.  The delivered graph
module stores the parent forest; it does not implement arbitrary group
evaluation or a mapping-class-group engine.  Keeping that distinction
prevents an algebraic compression from being mistaken for a solved
three-dimensional marking problem.

\subsection{Why simultaneous disjoint homology bases are insufficient}

The word ``multisurface'' in the announced architecture cannot be replaced
by ``a disjoint collection forming a full relative homology basis.''  The
obstruction is already visible in a product.

\begin{proposition}[A disjointness obstruction]
\label{prop:isotropic-obstruction}
Let $M=\Sigma_h\times I$, where $\Sigma_h$ is a closed connected orientable
surface of genus $h\geq1$.  Suppose pairwise disjoint, compact, properly
embedded, oriented surfaces $S_1,\ldots,S_k$ represent linearly independent classes
in $H_2(M,\partial M;\Q)$.  Then $k\leq h$, although that relative
homology group has dimension $2h$.
\end{proposition}
\begin{proof}
Poincar\'e--Lefschetz duality identifies their classes with linearly
independent classes $\alpha_i\in H^1(M;\Q)$.  Disjoint representatives
have zero mutual intersection, so $\alpha_i\smile\alpha_j=0$ for
$i\ne j$.  For $i=j$, the same vanishing follows from graded
commutativity over $\Q$.  Pullback along the inclusion of
$\Sigma_h\times\{1/2\}$ identifies $H^1(M;\Q)$ with its
$2h$-dimensional surface cohomology and preserves cup products.
The cup-product pairing, evaluated on $[\Sigma_h]$, is nondegenerate
and alternating.  An isotropic subspace of a $2h$-dimensional symplectic
space has dimension at most $h$: its inclusion in its symplectic
orthogonal gives $k\leq2h-k$.  This proves the result.
\end{proof}

The product in this example also embeds in $S^3$, as a small collar of an
embedded closed genus-$h$ surface.  Thus restricting attention to compact
submanifolds of $S^3$ does not remove the obstruction.
This classical obstruction identifies a concrete design requirement.
An implementation must support the intersections permitted by the
multisurface method, or justify a different hierarchy with equivalent
quantitative control.  Fast software for a disjoint special case does
not settle that requirement.

\subsection{The remaining logical gap}

The initial torus reducer starts with empty pattern and a verified
irreducible knot exterior.  A later cut normally introduces new boundary
data.  Erasing it when absorbing a ball may change the problem being
solved, even though the underlying unmarked manifold is unchanged.
Likewise, a cyclic word and a fibre flip do not automatically specify a
compatible system of boundary charts and all transition maps needed for
a gluing operation.  A nonorientable base with orientable total interval
bundle has no global fibre orientation; local frames and their transition
data are the appropriate objects.  The output contract of the native
compiler deliberately stops before these gluing assertions.

To close the global gap one must prove a complete surface-discovery
procedure with the required hierarchy length, establish bounds on the
pattern and its transport, and maintain those bounds through every
reconstruction and simplification.  The current contributions reduce
specific costs and remove specific missing source interfaces.  They
provide no shortcut around these remaining theorems.
